\section{Finite-width energy and the actual posterior}
\label{app:posterior}

The preceding sharp atomic and calibration theorems identify every
minimum-cost score. We now transfer their strict cost gaps to every
sufficiently large fixed width and compare the actual Gaussian
posterior integrals. The Gaussian variances are those in Section~1.

\begin{lemma}[Integer allocation of a finite atom representation]
Let $F=\sum_{j=1}^J c_j\tanh(a_j\cdot\bar x)$ be a finite
representation with nonzero $a_j,c_j$. Put
\[
 B=\sum_j|c_j|\|a_j\|_2,\qquad
 p_j=\frac{|c_j|}{\sqrt d\,\|a_j\|_2},\qquad p=\sum_jp_j.
\]
If $p<1$, then for every sufficiently large integer $N$ there is an
exact width-$N$ representation of $F$ with Gaussian energy tending to
$2\sqrt d\,B$ as $N\to\infty$.
\end{lemma}
\begin{proof}
Choose $r_j=\lceil Np_j\rceil$. For $N(1-p)\ge J$,
$\sum_jr_j\le N$. Use $r_j$ copies of the incoming vector $a_j$,
each with readout $c_j/r_j$, and pad with zero neurons. The output is
exactly $F$. Direct calculation gives
\[
 K_N-2\sqrt d\,B
 =\sum_j\frac{d\|a_j\|_2^2}{N}
                \frac{(r_j-Np_j)^2}{r_j}
 \le\frac dN\sum_j\|a_j\|_2^2.
\]
The right side tends to zero. The same statement applies after any
fixed small multiplicative increase of all readouts, provided the
resulting allocation sum remains less than one.
\end{proof}

For the balanced parity-orbit representation of a code score
$bg(z)=\sum_s\widehat g(s)\chi_{A_s}$, let $q_n$ denote its
positive parity correlation at the optimal scale $a_n$. The finite
orbit identity represents each parity using incoming vectors of norm
$a_n\sqrt n$ and total absolute readout $1/q_n$. Therefore
\[
 B=\sum_s w_{n_s}|\widehat g(s)|,
 \qquad
 p=\frac1{\sqrt d}\sum_s
       \frac{|\widehat g(s)|}{a_{n_s}\sqrt{n_s}\,q_{n_s}}.
\]
The balanced scalar asymptotics imply
$\inf_{n\ge n_*}a_n\sqrt n\,q_n>0$. Every baseline or rectangle
score in the construction has
$\sum_s|\widehat g(s)|=3/2$. Thus $p=O(d^{-1/2})<1/2$ for all
sufficiently large construction parameters. Integer allocation applies
uniformly to the finitely many rectangle witnesses needed for a fixed
instance, and also to a small strict-margin dilation of each witness.

Here is the output-space gap needed for the Bayesian argument. Fix a
history with at least one live point and pairwise distinct effective
rectangle costs. Choose a polynomial cap $B_{\rm cap}$ strictly above
all rectangle costs. Consider all pairs $(B,F)$ obeying the closed
atomic score inequalities of Appendix~\ref{app:physical}, the cap,
$\|F\|_\infty\le\sqrt{d+1}B_{\rm cap}$, the closed global margins,
and all revealed labels. This is a compact subset of a finite-dimensional
space. The calibration proof says that its unique
minimum-cost score is the cheapest live rectangle, of cost $C_*$.
That score is wrong at its exceptional live fiber.

Intersect the compact set with the additional condition that the score
is correct at the next chosen physical input in that fiber. This closed
subset has minimum cost $C_{\rm good}>C_*$. If it is empty, set
$C_{\rm good}=+\infty$. Actual feasible network representations of
cost at most $B_{\rm cap}$ belong to the compact set by the atomic
theorem and their full-cube and history margins. Representations above the cap have even
larger energy because $K_N\ge2\sqrt d\,B$. Consequently the energy of
every feasible network correct at the designated input is at least
\[
 E_{\rm good}=2\sqrt d\,\min(C_{\rm good},B_{\rm cap})
             >2\sqrt d\,C_*.
\]
Choose a small fixed dilation of the cheapest rectangle so that its
margins are strict and its variation cost remains below
$E_{\rm good}/(2\sqrt d)$. Integer allocation gives, for every
sufficiently large $N$, a strict feasible wrong network of energy
strictly below $E_{\rm good}$.

For completeness, fix one such $N$ and write $P=N(d+2)$ for the
parameter dimension. A small open ball around the strict feasible
witness lies in the conditioned version space and has energy bounded
by some $a<E_{\rm good}$. The correct region has energy at least
$E_{\rm good}$. Choose $\rho\in(0,1)$ so that
$(1-\rho)E_{\rm good}>a$. Its unnormalized Gaussian integral is at most
\[
 e^{-(1-\rho)E_{\rm good}/(2\sigma^2)}
 \int_{\mathbb R^P}e^{-\rho K_N(\theta)/(2\sigma^2)}\,d\theta
 =C_{N,\rho}\sigma^P
    e^{-(1-\rho)E_{\rm good}/(2\sigma^2)},
\]
whereas the version-space integral is at least the ball volume times
$e^{-a/(2\sigma^2)}$. Their ratio tends to zero. Thus the actual
conditioned Gaussian posterior predicts the wrong sign with probability
tending to one as $\sigma\downarrow0$ at this fixed width.

There are finitely many logical rounds in each fixed construction.
Take the maximum of their width thresholds. The posterior conclusion
then holds for every fixed width beyond that threshold and every round.


Pairwise distinctness of every candidate is stronger than necessary.
If there are $T$ distinct effective cost levels, first reveal every
candidate fiber except one representative per level. At each subsequent
round, take one physical input in the cheapest surviving representative
fiber, then reveal the remainder of that fiber before continuing. All
these labels are target labels, all inputs are distinct, and the
same analysis applies at each designated prediction. This gives
$T$ designated mistakes with posterior probability tending to one at
each round. Section~\ref{sec:cycles} establishes exponentially many such levels
in the logical dimension for balanced tanh costs on its path column code.
